\section*{Chapitre \ref{ch:g3:flowers-fruits-seeds} --- Fleurs, fruits et graines}

\begin{solution}{exo:g3:flowers-fruits-seeds:1}
Les \omterm{def:g3:flowers-fruits-seeds:parts}{pétales} (colorés et parfumés, ils attirent les \omterm{prop:g3:sorting-living-things:groups}{insectes})~; les
\omterm{def:g3:flowers-fruits-seeds:parts}{étamines} (fines \omterm{def:g1:plants-around-us:parts}{tiges} qui portent le \omterm{def:g3:flowers-fruits-seeds:parts}{pollen})~; le \omterm{def:g3:flowers-fruits-seeds:parts}{pistil} au centre
(la petite fiole qui contient les futures \omterm{prop:g3:flowers-fruits-seeds:transform}{graines}).
\end{solution}

\begin{solution}{exo:g3:flowers-fruits-seeds:2}
Le \omterm{def:g3:flowers-fruits-seeds:parts}{pollen} est la fine poussière jaune portée par les \omterm{def:g3:flowers-fruits-seeds:parts}{étamines}. La
pollinisation, c'est le \omterm{def:g3:flowers-fruits-seeds:parts}{pollen} qui atteint le \omterm{def:g3:flowers-fruits-seeds:parts}{pistil} --- livré le
plus souvent par les \omterm{prop:g3:sorting-living-things:groups}{insectes} de passage, parfois par le vent.
\end{solution}

\begin{solution}{exo:g3:flowers-fruits-seeds:3}
(a) Ils se fanent et tombent. (b) Il gonfle et devient le \omterm{prop:g3:flowers-fruits-seeds:transform}{fruit}.
(c) Elles deviennent les \omterm{prop:g3:flowers-fruits-seeds:transform}{graines} à l'intérieur du \omterm{prop:g3:flowers-fruits-seeds:transform}{fruit}.
\end{solution}

\begin{solution}{exo:g3:flowers-fruits-seeds:4}
Avril~: \omterm{def:g3:flowers-fruits-seeds:parts}{fleurs} blanches, visites d'abeilles --- la pollinisation.
Mai~: les \omterm{def:g3:flowers-fruits-seeds:parts}{pétales} tombent~; le \omterm{def:g3:flowers-fruits-seeds:parts}{pistil} de chaque \omterm{def:g3:flowers-fruits-seeds:parts}{fleur} pollinisée
gonfle en une petite boule verte. Juin~: la boule est devenue une
cerise rouge et sucrée, avec le noyau --- la \omterm{def:g2:from-seed-to-plant:seed}{graine} --- dedans.
\end{solution}

\begin{solution}{exo:g3:flowers-fruits-seeds:5}
Parce que sans visites d'\omterm{prop:g3:sorting-living-things:groups}{insectes} la plupart des \omterm{def:g3:flowers-fruits-seeds:parts}{fleurs} ne sont
jamais pollinisées, et qu'une \omterm{def:g3:flowers-fruits-seeds:parts}{fleur} non pollinisée tombe sans
devenir un \omterm{prop:g3:flowers-fruits-seeds:transform}{fruit}. Ce sont les livraisons de \omterm{def:g3:flowers-fruits-seeds:parts}{pollen} des abeilles qui
changent les \omterm{def:g3:flowers-fruits-seeds:parts}{fleurs} en cerises.
\end{solution}

\begin{solution}{exo:g3:flowers-fruits-seeds:6}
Coupe-le et cherche des \omterm{prop:g3:flowers-fruits-seeds:transform}{graines}~: des \omterm{prop:g3:flowers-fruits-seeds:transform}{graines} dedans veut dire qu'il
a poussé à partir d'un \omterm{def:g3:flowers-fruits-seeds:parts}{pistil} --- un \omterm{prop:g3:flowers-fruits-seeds:transform}{fruit}~; aucune veut dire une
autre partie de la \omterm{def:g1:plants-around-us:plant}{plante}. La tomate est pleine de \omterm{prop:g3:flowers-fruits-seeds:transform}{graines}~: un
\omterm{prop:g3:flowers-fruits-seeds:transform}{fruit}. La carotte n'en a pas~: pas un \omterm{prop:g3:flowers-fruits-seeds:transform}{fruit}, mais une \omterm{def:g1:plants-around-us:parts}{racine}.
\end{solution}

\begin{solution}{exo:g3:flowers-fruits-seeds:7}
Concombre~: des \omterm{prop:g3:flowers-fruits-seeds:transform}{graines} --- un \omterm{prop:g3:flowers-fruits-seeds:transform}{fruit}. Laitue~: pas de \omterm{prop:g3:flowers-fruits-seeds:transform}{graines} ---
des \omterm{def:g1:plants-around-us:parts}{feuilles}, pas un \omterm{prop:g3:flowers-fruits-seeds:transform}{fruit}. Gousse de haricot vert~: une rangée de
\omterm{prop:g3:flowers-fruits-seeds:transform}{graines} --- un \omterm{prop:g3:flowers-fruits-seeds:transform}{fruit}. Pomme de terre~: pas de \omterm{prop:g3:flowers-fruits-seeds:transform}{graines} --- une \omterm{def:g1:plants-around-us:parts}{tige}
souterraine renflée, pas un \omterm{prop:g3:flowers-fruits-seeds:transform}{fruit}.
\end{solution}

\begin{solution}{exo:g3:flowers-fruits-seeds:8}
Réponse libre. Les surprises habituelles~: la tomate, le concombre
et la gousse de haricot qui se révèlent être des \omterm{prop:g3:flowers-fruits-seeds:transform}{fruits}~; la fraise
qui porte ses minuscules \omterm{prop:g3:flowers-fruits-seeds:transform}{graines} à l'extérieur.
\end{solution}

\begin{solution}{exo:g3:flowers-fruits-seeds:9}
Presque pas de prunes. Un \omterm{prop:g3:flowers-fruits-seeds:transform}{fruit} ne peut venir que du \omterm{def:g3:flowers-fruits-seeds:parts}{pistil} d'une
\omterm{def:g3:flowers-fruits-seeds:parts}{fleur} pollinisée~; les \omterm{def:g3:flowers-fruits-seeds:parts}{fleurs} mortes en avril, il ne reste rien à
transformer, si bel été soit-il.
\end{solution}

\begin{solution}{exo:g3:flowers-fruits-seeds:10}
Pour l'arbre, la pomme est une façon de fabriquer et de protéger des
\omterm{prop:g3:flowers-fruits-seeds:transform}{graines}. En son cœur se trouvent les pépins --- des \omterm{prop:g3:flowers-fruits-seeds:transform}{graines} --- et
chaque pépin, planté, pourrait \omterm{def:g2:from-seed-to-plant:germination}{germer} en un nouveau pommier. Que
nous trouvions l'emballage délicieux est notre chance (et, pour
l'arbre, une façon de faire emporter ses \omterm{prop:g3:flowers-fruits-seeds:transform}{graines} loin de chez lui).
\end{solution}
